%!TEX root = ../main.tex
\section{Introduction}
\label{sec:introduction}
%Given the increasing use of multimodal AI systems in scientific research~\citep{anthropicclauderesearch2026,google2024gemini_deepresearch}, it has become imperative to understand how well vision-language models perform on scientific figures and to characterise their failure modes. This is critical because scientific figures often encode the central quantitative claims of a paper, and errors in their interpretation can propagate into inaccurate downstream summaries, comparisons, and research conclusions~\citep{wang2024charxiv,huang2024scifibench,tang2025chartmuseum}. A growing body of benchmarks addresses this need across chart question answering, scientific chart understanding, visual mathematical reasoning, and college-level multimodal reasoning~\citep{masry2022chartqa,lu2024mathvista,yue2024mmmu,wang2024charxiv,tang2025chartmuseum}. These benchmarks have substantially advanced evaluation, but are often closed-form and centred on aggregate accuracy.

The growing use of multimodal AI systems in scientific research~\citep{yan2025multimodalscience,zhang2025scientificllms,Eger2025TransformingSW,greisinger2026tikzilla,zhang2025scimage} has increased the need to evaluate how vision-language models (VLMs) interpret scientific figures and where they fail. Scientific figures often contain the core quantitative evidence of a paper, making errors in interpretation especially consequential for downstream summarisation, comparison, and scientific reasoning~\citep{wang2024charxiv,huang2024scifibench,tang2025chartmuseum}. Recent benchmarks have advanced evaluation across chart understanding, visual reasoning, and multimodal scientific QA~\citep{masry2022chartqa,lu2024mathvista,yue2024mmmu,wang2024charxiv,tang2025chartmuseum}, but most focus on closed-form tasks and aggregate accuracy.

Accuracy under clean conditions, however, does not capture how models behave when evidence is incomplete or misleading. Figure~\ref{fig:hook} illustrates this gap. In Figure~\ref{fig:hook-original},
the label ``Academic Funding'' is clearly visible whereas in Figure~\ref{fig:hook-blur} 
the single label is selectively blurred. \po{GPT-5.2 correctly describes and reasons about the original figure but, when asked about the blurred region, answers ``Customer Support'', a category absent from the chart, without acknowledging uncertainty or unreadable evidence.} 

%As depicted in figure~\ref{fig:hook}, accuracy under clean conditions may not be a sufficient metric for evaluating these models. (a), the category ``Academic Funding'' is legible; in the degraded variant (b), that single label has been selectively blurred. Given the clean chart, GPT-5.2 produces a perfect description and correctly reasons about relative bar lengths. Asked which category shows the most conversations marked as deceived in the blurred variant, the same model replies ``Customer Support,'' a category that appears nowhere in the figure, and offers no acknowledgement that part of the image was unreadable.

%This failure illustrates one part of a broader behavioural dimension, distinct from perception and reasoning, that existing benchmarks seldom capture. We operationalise this dimension through what we term the Admittance--Resistance--Inductance (A-R-I) framework (\S\ref{sec:ari}), which measures a model's epistemic honesty under visual uncertainty (Admittance), its robustness to false premises and misleading context (Resistance), and its capacity for contextual inference when visual information is degraded (Inductance). These behaviours matter in research workflows because local perceptual uncertainty can propagate into unsupported downstream conclusions, a concern echoed by deployed vision systems in which high performance under expected conditions has coexisted with failures under degraded or unusual visual inputs~\citep{beede2020human,nhtsa2024teslafsd}.

This failure reflects a dimension distinct from perception and reasoning that existing benchmarks rarely measure directly. We formalise this dimension through the Admittance--Resistance--Inductance (A-R-I) framework (\S\ref{sec:ari}), which evaluates whether models acknowledge insufficient evidence (Admittance), resist misleading context or false premises (Resistance), and its capacity for contextual inference when visual information is degraded (Inductance). \po{While we instantiate this framework in the context of scientific figures, these capabilities are broadly applicable to multimodal settings in which models must distinguish observed evidence from uncertainty and inference, including but not limited to healthcare (\textit{medical imaging and diagnosis}), autonomous systems (\textit{scene understanding under adverse conditions}), document intelligence (\textit{OCR and form understanding}), education (\textit{diagram and textbook interpretation}), remote sensing (\textit{satellite image analysis}), robotics (\textit{reasoning under partial observations}), and legal or financial document analysis, where local perceptual errors can propagate into confident but unsupported conclusions~\citep{beede2020human,nhtsa2024teslafsd}.}

%These properties are especially important in scientific workflows, where local perceptual errors can propagate into unsupported conclusions ~\citep{beede2020human,nhtsa2024teslafsd}. 
%\todo{SE: I think this is interesting. The framework could apply more broadly, beyond science - ANANYA}

%Across 8 frontier models and more than 23,000 evaluation instances constructed from 250 scientific figures, we find that models with similar perception and reasoning scores, measured through open-ended descriptions and figure-based questions, diverge sharply on behaviour (as measured by our proposed A-R-I framework). The highest-scoring model on description quality fabricates answers for elements it cannot see 98\% of the time, while a comparably-scoring model admits uncertainty 90\% of the time.

Using 250 scientific figures with over 34,000 evaluation setups  across eight frontier models, we show that models with similar perception and reasoning performance can behave very differently under uncertainty. The top-performing model on description quality fabricates answers for unreadable content in 96\% of cases, whereas a similarly performing model admits uncertainty 71\% of the time. Our contributions are as follows:

% \begin{enumerate}[nosep]
%     \item \textsc{SciFig-Eval}, a benchmark for scientific figure understanding that evaluates VLMs across perception, reasoning, and behaviour. The benchmark comprises 250 annotated scientific figures, open-ended figure descriptions scored with MQM-adapted criteria, and 1,000 figure-based capability questions spanning counting, computation, comparison, and pattern analysis.

%     \item A controlled stress-test suite for scientific figures, covering 1,243 transformed or in-paper figure cases, 100 caption-bias cases, and 750 resistance probes over contradictory, non-existent, and unanswerable premises. The transformed cases include low-contrast, noise, and rotation variants, together with figure-on-page and figure-blurred-on-page conditions that test whether models rely on surrounding document context when the figure is degraded.

%     \item The Admittance--Resistance--Inductance (A-R-I) behavioural framework (\S\ref{sec:ari}), a three-axis model that decomposes behaviour into epistemic honesty under visual uncertainty, robustness to misleading context, and bounded inference from partial evidence. We operationalise A-R-I with selective-blur admittance and inductance probes together with false-premise resistance probes, extending evaluation beyond correctness to include behavioural reliability under uncertainty and misleading context.
% \end{enumerate}

\begin{enumerate}[label=(\roman*)]
    \item We introduce \textsc{SciFig-Eval}\footnote{\url{https://mango-dune-05852c00f.7.azurestaticapps.net/\#dataset}}, a benchmark for scientific figure understanding to evaluate VLMs across perception, reasoning, and behavioural reliability. It contains 250 annotated figures, MQM-based open-ended descriptions, and 1,000 figure-grounded reasoning questions covering counting, computation, comparison, and pattern analysis.

    \item We develop a controlled stress-test 
    (\S\ref{sec:mqm})
    suite with transformed figures, caption-bias settings, and false-premise probes to evaluate robustness under degraded visual evidence.

    \item We propose the Admittance-Resistance-Inductance (A-R-I) framework (\S\ref{sec:ari}), by decomposing behaviour into uncertainty acknowledgment, resistance to misleading context and inference from partial evidence through selective-blur and false-premise probes.
\end{enumerate}
